\documentclass[11pt]{article}
\usepackage{ap-lit}
\title{Unit 2: Poetry: Question Set}
\author{Paul Wang}
\date{November 2026}
\docpart{Unit 2 \textperiodcentered\ Question Set}

\begin{document}
\maketitle

\begin{abstract}
Three exam-style multiple-choice sets on public-domain poems, followed by a full Poetry Analysis essay. Poem 1 is Shakespeare's Sonnet 73, a tightly built argument in three images; Poem 2 is Thomas Hardy's ``The Darkling Thrush,'' a poem that turns on a single sound; Poem 3 is Paul Laurence Dunbar's ``We Wear the Mask,'' a rondeau built on a refrain. The questions follow the exam's emphasis on the speaker, shifts, and what comparisons transfer. Allow about 27 minutes for the multiple choice and 40 minutes for the essay. Line numbers are printed every five lines and are cited in the questions.
\end{abstract}

\scoringline{MCQ: \blank\ / \qtotal \qquad Essay: \blank\ / 6 \qquad Target: 18+ MCQ, 4+ essay}

\section{Poem 1: William Shakespeare, Sonnet 73}

\begin{poem}{Questions 1--9.\quad Read the poem carefully before you choose your answers.}{The following poem is Sonnet 73 from William Shakespeare's sonnet sequence, published in 1609.}
\pl{That time of year thou mayst in me behold}
\pl{When yellow leaves, or none, or few, do hang}
\pl{Upon those boughs which shake against the cold,}
\pl{Bare ruin'd choirs, where late the sweet birds sang.}
\pl{In me thou see'st the twilight of such day}
\pl{As after sunset fadeth in the west;}
\pl{Which by and by black night doth take away,}
\pl{Death's second self, that seals up all in rest.}
\pl{In me thou see'st the glowing of such fire,}
\pl{That on the ashes of his youth doth lie,}
\pl{As the death-bed, whereon it must expire,}
\pl{Consum'd with that which it was nourish'd by.}
\pli{This thou perceiv'st, which makes thy love more strong,}
\pli{To love that well, which thou must leave ere long.}
\end{poem}

\begin{q}{2}{2.1 Speaker and listener}{``Thou'' who sees him and whose ``love'' is named in line 13: a beloved who will lose him.}
The speaker addresses the poem to
\ch{Death, whom he asks for more time}{a listener who loves him and will soon lose him}{autumn, personified as a ruined church}{his own younger self}
\end{q}

\begin{q}{1}{2.2 Structure}{Year, then day, then a dying fire: each quatrain is a new image of decline, each on a shorter time scale.}
The three quatrains (lines 1--4, 5--8, and 9--12) are best described as
\ch{three images of the speaker's decline, each measured on a shorter span of time}{three separate stories from the speaker's past}{three addresses to three different listeners}{three contradictory arguments the couplet rejects}
\end{q}

\begin{q}{4}{2.4 Word choice}{The speaker hesitates and revises (``or none, or few''), as if reluctant to measure the loss.}
The phrase ``yellow leaves, or none, or few'' (line 2) suggests a speaker who is
\ch{listing the seasons in their natural order}{certain that all his leaves are gone}{celebrating the colors of autumn}{hesitating, as if unwilling to state plainly how far his decline has gone}
\end{q}

\begin{q}{3}{2.5 Metaphor}{Bare branches = ruined choir stalls where song once was.}
In line 4, ``Bare ruin'd choirs, where late the sweet birds sang'' compares
\ch{the speaker's voice to a choir of birds}{monks to migrating birds}{bare branches to the ruined choir of a church, once full of song}{the listener to a singer who has left}
\end{q}

\begin{q}{1}{2.5 Metaphor}{Night is death's double: it ``seals up all in rest.''}
``Death's second self'' (line 8) refers to
\chh{black night}{the sunset}{the speaker}{the listener}
\end{q}

\begin{q}{2}{2.5 Paradox}{The ashes of what fed the fire now smother it: life's fuel becomes what ends it.}
Line 12, ``Consum'd with that which it was nourish'd by,'' presents the speaker's aging as
\ch{a sudden accident}{a paradox in which what once sustained his life now extinguishes it}{a punishment for his past}{a process that can be reversed}
\end{q}

\begin{q}{4}{2.2 Repetition}{``In me thou see'st'' repeated at lines 5 and 9 insists that the listener look at him.}
The repetition of ``In me thou see'st'' (lines 5 and 9) chiefly serves to
\ch{introduce a new speaker in each quatrain}{create an internal rhyme with ``sweet''}{shift the poem into the past tense}{insist that the listener look at the speaker's decline in image after image}
\end{q}

\begin{q}{3}{2.2 Couplet}{The couplet turns from his decline to what it does to the listener's love.}
The couplet (lines 13--14) shifts the focus from
\ch{the listener's youth to the speaker's age}{fire to water}{the speaker's decline to its effect on the listener's love}{praise to complaint}
\end{q}

\begin{q}{1}{2.6 Literary argument}{The decline becomes his argument: love grows stronger because it must soon let go.}
Which interpretation of the poem is best supported?
\ch{The speaker turns the evidence of his own decline into a reason the listener's love should grow stronger.}{The speaker regrets that the listener has stopped loving him.}{The speaker believes his youth will return in spring.}{The speaker urges the listener to leave him before he dies.}
\end{q}

\section{Poem 2: Thomas Hardy, ``The Darkling Thrush''}

\begin{poem}{Questions 10--18.\quad Read the poem carefully before you choose your answers.}{The following poem by Thomas Hardy is dated December 1900, the last month of the nineteenth century. A ``coppice'' is a thicket of small trees; ``bine-stems'' are the stems of climbing plants.}
\pl{I leant upon a coppice gate}
\pli{When Frost was spectre-gray,}
\pl{And Winter's dregs made desolate}
\pli{The weakening eye of day.}
\pl{The tangled bine-stems scored the sky}
\pli{Like strings from broken lyres,}
\pl{And all mankind that haunted nigh}
\pli{Had sought their household fires.}
\stanza
\pl{The land's sharp features seemed to be}
\pli{The Century's corpse outleant,}
\pl{His crypt the cloudy canopy,}
\pli{The wind his death-lament.}
\pl{The ancient pulse of germ and birth}
\pli{Was shrunken hard and dry,}
\pl{And every spirit upon earth}
\pli{Seemed fervourless as I.}
\stanza
\pl{At once a voice outburst among}
\pli{The bleak twigs overhead}
\pl{In a full-hearted evensong}
\pli{Of joy illimited;}
\pl{An aged thrush, frail, gaunt, and small,}
\pli{In blast-beruffled plume,}
\pl{Had chosen thus to fling his soul}
\pli{Upon the growing gloom.}
\stanza
\pl{So little cause for carollings}
\pli{Of such ecstatic sound}
\pl{Was written on terrestrial things}
\pli{Afar or nigh around,}
\pl{That I could think there trembled through}
\pli{His happy good-night air}
\pl{Some blessed Hope, whereof he knew}
\pli{And I was unaware.}
\src{December 1900}
\end{poem}

\begin{q}{4}{2.1 Situation}{Alone at a gate at dusk in winter, after everyone has gone indoors.}
In the first stanza, the speaker is
\ch{walking home to join his family}{watching a crowd gather for a celebration}{describing a scene he has only imagined}{alone outdoors at a winter dusk after others have gone inside}
\end{q}

\begin{q}{2}{2.5 Simile}{Stems like the strings of broken lyres: music and art now broken and silent.}
The simile in lines 5--6, comparing the bine-stems to ``strings from broken lyres,'' suggests
\ch{that the speaker is a skilled musician}{a world whose music and poetry have broken down into silence}{the sound of wind playing through the trees}{the beauty of a winter evening}
\end{q}

\begin{q}{3}{2.5 Extended metaphor}{The landscape becomes the century's corpse, clouds its crypt, wind its lament.}
In lines 9--12, the landscape is presented as
\ch{a garden in early spring}{a battlefield after a war}{the corpse of the dying century, laid out in a crypt and mourned by the wind}{a church where the speaker prays}
\end{q}

\begin{q}{1}{2.1 Speaker}{``Fervourless as I'': the speaker measures every spirit by his own lack of hope.}
The phrase ``every spirit upon earth / Seemed fervourless as I'' (lines 15--16) suggests that
\ch{the speaker sees his own lack of hope reflected in the world around him}{the speaker is more hopeful than anyone else}{the speaker has spoken with every person on earth}{the speaker is describing the thrush}
\end{q}

\begin{q}{4}{2.2 Shift}{``At once'' marks the turn from desolation to sudden song.}
The phrase ``At once'' (line 17) marks a shift from
\ch{hope to despair}{the speaker's past to his future}{a description of the thrush to a description of the sky}{a world of silence and death to a sudden outburst of song}
\end{q}

\begin{q}{2}{2.4 Imagery}{The thrush is as worn as the scene, so its joyful song is all the more surprising.}
The description of the thrush as ``aged \ldots\ frail, gaunt, and small, / In blast-beruffled plume'' (lines 21--22) chiefly emphasizes that
\ch{the bird is young and newly hatched}{the singer shares the bleakness of the scene, which makes its joy more striking}{the speaker pities the bird and wants to help it}{the bird is a symbol of the dying century}
\end{q}

\begin{q}{3}{2.4 Word choice}{``Fling'' is reckless and generous: a defiant gesture against the gloom.}
In line 23, the verb ``fling'' suggests that the thrush's song is
\ch{carefully rehearsed}{reluctant and weak}{a reckless, generous act of defiance against the darkness}{a warning to other birds}
\end{q}

\begin{q}{1}{2.1 Tone of the ending}{``I could think'' and ``I was unaware'': hope allowed as possible, not claimed.}
The final stanza is best described as
\ch{tentative: the speaker allows that hope may exist while remaining outside it himself}{triumphant: the speaker is converted to the thrush's joy}{mocking: the speaker ridicules the bird's foolish song}{despairing: the speaker denies any possibility of hope}
\end{q}

\begin{q}{4}{2.4 Diction}{Funeral words in stanza 2 give way to song and celebration words in stanzas 3--4.}
The words ``corpse,'' ``crypt,'' and ``death-lament'' in stanza 2 contrast most directly with which group of words from stanzas 3 and 4?
\ch{``bleak,'' ``gaunt,'' ``gloom''}{``coppice,'' ``bine-stems,'' ``canopy''}{``Frost,'' ``Winter's,'' ``cold''}{``evensong,'' ``carollings,'' ``ecstatic''}
\end{q}

\section{Poem 3: Paul Laurence Dunbar, ``We Wear the Mask''}

\begin{poem}{Questions 19--25.\quad Read the poem carefully before you choose your answers.}{The following poem by Paul Laurence Dunbar (1872--1906), one of the first African American poets to win a national readership, was published in 1895.}
\pl{We wear the mask that grins and lies,}
\pl{It hides our cheeks and shades our eyes,---}
\pl{This debt we pay to human guile;}
\pl{With torn and bleeding hearts we smile,}
\pl{And mouth with myriad subtleties.}
\stanza
\pl{Why should the world be over-wise,}
\pl{In counting all our tears and sighs?}
\pl{Nay, let them only see us, while}
\pli{\hspace*{1.5em}We wear the mask.}
\stanza
\pl{We smile, but, O great Christ, our cries}
\pl{To thee from tortured souls arise.}
\pl{We sing, but oh the clay is vile}
\pl{Beneath our feet, and long the mile;}
\pl{But let the world dream otherwise,}
\pli{\hspace*{1.5em}We wear the mask!}
\end{poem}

\begin{q}{2}{2.1 Speaker}{``We'' speaks for a group forced to hide suffering behind a public face.}
The speaker of the poem speaks
\ch{as a single actor preparing for a play}{on behalf of a group who must hide their suffering behind a public face}{to a friend who has betrayed him}{as the world that watches the mask-wearers}
\end{q}

\begin{q}{3}{2.5 Personification}{The mask itself ``grins and lies'': the public face acts as a deceiver separate from the self.}
In line 1, describing the mask as one ``that grins and lies'' suggests that
\ch{the speaker enjoys deceiving the world}{the mask is a literal theatrical prop}{the public face acts as a deceiver separate from the self behind it}{the world is fooled only briefly}
\end{q}

\begin{q}{1}{2.4 Word choice}{A ``debt'' to ``human guile'': an obligation exacted by others' deceit and cruelty.}
The phrase ``This debt we pay to human guile'' (line 3) presents wearing the mask as
\ch{an obligation forced on the speakers by the deceit and cruelty of others}{a voluntary gift the speakers enjoy giving}{a financial arrangement between neighbors}{a punishment for the speakers' own lies}
\end{q}

\begin{q}{4}{2.1 Tone}{The rhetorical question refuses the world the sight of their pain: bitter defiance.}
The rhetorical question in lines 6--7 conveys a tone of
\ch{curious uncertainty}{cheerful acceptance}{humble apology}{bitter defiance, refusing the world the sight of their pain}
\end{q}

\begin{q}{2}{2.2 Shift}{The third stanza turns from the world to Christ, the one witness to their cries.}
Lines 10--11 (``We smile, but, O great Christ, our cries / To thee from tortured souls arise'') mark a shift from
\ch{grief to joy}{addressing the watching world to calling on Christ, the one witness to the suffering behind the mask}{the present to the distant past}{the plural ``we'' to a single ``I''}
\end{q}

\begin{q}{3}{2.2 Refrain}{The short refrain returns; the final exclamation turns statement into declaration.}
The refrain ``We wear the mask'' ends with a period in line 9 and an exclamation point in line 15. The change chiefly suggests that the speaker's final use of the phrase is
\ch{a question the speaker cannot answer}{an admission that the mask has been removed}{a more forceful, defiant declaration}{a quotation from the watching world}
\end{q}

\begin{q}{4}{2.4 Imagery}{Vile clay underfoot and a long mile: a long, degrading road of suffering.}
The image ``the clay is vile / Beneath our feet, and long the mile'' (lines 12--13) suggests
\ch{a pleasant journey home}{the speakers' love of the countryside}{the speakers' wish to travel abroad}{a long and degrading road of suffering}
\end{q}

\section{Free Response: Poetry Analysis}

\begin{frq}{Question 1 \textperiodcentered\ Poetry Analysis \textperiodcentered\ 6 points \textperiodcentered\ 40 minutes}
In the following sonnet by Christina Rossetti, published in 1862, the speaker addresses someone she loves. Read the poem carefully. Then, in a well-written essay, analyze how Rossetti uses poetic elements and techniques to convey the speaker's complex attitude toward being remembered.

In your response you should do the following:
\begin{itemize}[leftmargin=1.4em]
\item Respond to the prompt with a thesis that presents a defensible interpretation.
\item Select and use evidence to support your line of reasoning.
\item Explain how the evidence supports your line of reasoning.
\item Use appropriate grammar and punctuation in communicating your argument.
\end{itemize}
\end{frq}

\begin{poem}{Christina Rossetti, ``Remember'' (1862)}{}
\pl{Remember me when I am gone away,}
\pli{Gone far away into the silent land;}
\pli{When you can no more hold me by the hand,}
\pl{Nor I half turn to go yet turning stay.}
\pl{Remember me when no more day by day}
\pli{You tell me of our future that you planned:}
\pli{Only remember me; you understand}
\pl{It will be late to counsel then or pray.}
\pl{Yet if you should forget me for a while}
\pli{And afterwards remember, do not grieve:}
\pli{For if the darkness and corruption leave}
\pli{A vestige of the thoughts that once I had,}
\pl{Better by far you should forget and smile}
\pli{Than that you should remember and be sad.}
\end{poem}

\makekey

\subsection*{Essay scoring guide (6 points)}
\small
\begin{compare}{Row}{Earns the point(s)}{Does not earn the point(s)}
A. Thesis (1) & A defensible interpretation of the \emph{complex} attitude, e.g., that the speaker begins by demanding to be remembered and ends by releasing the beloved from that demand, so that her love finally values the beloved's happiness over her own survival in memory. & ``Rossetti uses repetition, imagery, and a volta to show the speaker's attitude''; ``The speaker wants to be remembered''; a summary of the poem. \\
B. Evidence and commentary (4) & \textbf{4}: specific evidence from both the octave and the sestet (the insistent imperative ``Remember me'' and its repetitions; the euphemism ``the silent land''; the tender gesture of line 4, ``half turn to go yet turning stay''; the lost future ``that you planned''; the volta at ``Yet''; ``darkness and corruption''; the closing antithesis of ``forget and smile'' against ``remember and be sad'') with commentary that consistently explains how form, repetition, imagery, and the turn build the interpretation. \textbf{3}: the same evidence, but only some is explained. \textbf{2}: some specific evidence; commentary is uneven or partly mistaken (e.g., treating the sestet as a continuation of the octave). \textbf{1}: paraphrase with little explanation. & Identifying the sonnet form or the volta without explaining what the turn does; listing repetitions without connecting them to the attitude. \\
C. Sophistication (1) & Develops a tension across the essay, e.g., that the sestet's generosity is still a way of being remembered (asking the beloved to ``forget'' keeps her in mind); that the Petrarchan structure itself enacts the change of heart; or that the poem's love is defined by giving up the very thing the octave asks for. & A concluding sentence that ``love is selfless''; a single unexplained mention of the Petrarchan form. \\
\end{compare}

\begin{modelessay}[A 6-point response, in outline]
\textbf{Thesis.} Rossetti's sonnet turns on its own opening command: the octave's insistent ``Remember me'' clings to the beloved, while the sestet, pivoting on ``Yet,'' releases him, so that the speaker's love proves itself by surrendering the remembrance it first demanded.

\textbf{Body 1: the octave holds on.} The imperative ``Remember me'' opens lines 1 and 5 and returns, softened, as ``Only remember me'' in line 7; the repetition is the speaker's grip. Even the euphemism ``the silent land'' avoids the word ``death,'' as if naming it would loosen her hold. Line 4, ``Nor I half turn to go yet turning stay,'' compresses their whole relationship into a gesture of reluctant parting, and its doubled ``turn'' anticipates the poem's own turn.

\textbf{Body 2: the volta lets go.} ``Yet'' at line 9 begins the sestet and reverses the octave's logic. Forgetting, which the octave dreaded, is now permitted; grief, which it implicitly asked for, is now forbidden: ``do not grieve.'' The blunt ``darkness and corruption'' replaces the gentle ``silent land,'' a sign that the speaker can now face death plainly because she is no longer bargaining with it.

\textbf{Body 3: the closing antithesis.} The final couplet balances ``forget and smile'' against ``remember and be sad,'' and chooses the first. The rhyme of ``while'' with ``smile'' links the beloved's brief forgetting to his happiness, and the speaker's last wish is for his good, not her memory.

\textbf{Why it earns Row C.} The essay keeps the tension alive: the sestet's release is itself an act of devotion that the beloved will remember, and the Petrarchan octave--sestet structure carries the speaker's change of heart.
\end{modelessay}

\end{document}
